\chapter{Réalisation des modules avancés}

% ============================================================
\section{Introduction}
% ============================================================

Ce chapitre présente la réalisation des sprints~4 et~5. Le Sprint~4 développe les
fonctionnalités de notifications temps réel, collaboration, dashboards et calendrier.
Le Sprint~5 implémente le moteur d'intelligence artificielle hybride de la plateforme :
recommandation d'assignation par Machine Learning, estimation de durée par heuristiques
enrichies, détection de risques par règles métier, et assistant conversationnel LLM.

% ============================================================
\section{Sprint 4 : Notifications, collaboration et communication}
% ============================================================

\subsection{Spécifications fonctionnelles}

\begin{table}[H]
  \centering
  \caption{Sprint Backlog — Sprint 4 (39 story points)}
  \label{tab:backlog_s4}
  \begin{tabularx}{\textwidth}{C{1.2cm}L{5.8cm}L{2.2cm}C{2cm}C{1.5cm}}
    \toprule
    \thead{ID} & \thead{User Story} & \thead{Acteur} & \thead{Priorité} & \thead{Points} \\
    \midrule
    \rowcolor{lightgray}
    US-17 & Recevoir des notifications in-app temps réel  & Tous   & Must   & 8 \\
    US-18 & Configurer ses préférences de notification    & Tous   & Should & 3 \\
    \rowcolor{lightgray}
    US-19 & Consulter le calendrier mensuel               & Member & Should & 5 \\
    US-20 & Accéder à un dashboard personnel              & Member & Should & 8 \\
    \rowcolor{lightgray}
    US-21 & Accéder au dashboard administrateur           & Admin  & Should & 5 \\
    US-22 & Utiliser la recherche globale cross-projet    & Tous   & Should & 5 \\
    \rowcolor{lightgray}
    US-23 & Consulter le rapport d'avancement (Client)   & Client & Should & 5 \\
    \midrule
    \multicolumn{4}{r}{\textbf{Total :}} & \textbf{39} \\
    \bottomrule
  \end{tabularx}
\end{table}

Les fonctionnalités complémentaires réalisées : suivi de tâches
(\texttt{WatcherController} — \texttt{POST /tasks/\{id\}/watch}),
suivi de projets (\texttt{FollowerController} — \texttt{POST /projects/\{id\}/follow}),
performance des membres (\texttt{SpecialtyController::membersPerformance}) et
analytics de charge de travail.

\subsection{Conception}

\begin{figure}[H]
  \centering
  \IfFileExists{figures/sprint4_use_case.png}{%
    \includegraphics[width=0.88\textwidth]{figures/sprint4_use_case.png}%
  }{%
    \fbox{\parbox{0.88\textwidth}{\centering\vspace{1.8cm}
      \textbf{[IMAGE À INSÉRER]}\\[0.4em]
      \texttt{figures/sprint4\_use\_case.png}\\[0.3em]
      \textit{Cas d'utilisation Sprint 4 : notifications, calendrier, dashboards,
      recherche, suivi tâches/projets}
    \vspace{1.8cm}}}%
  }
  \caption{Diagramme de cas d'utilisation — Sprint 4}
  \label{fig:uc_s4}
\end{figure}

\begin{figure}[H]
  \centering
  \IfFileExists{figures/sprint4_sequence_notif.png}{%
    \includegraphics[width=0.90\textwidth]{figures/sprint4_sequence_notif.png}%
  }{%
    \fbox{\parbox{0.90\textwidth}{\centering\vspace{1.8cm}
      \textbf{[IMAGE À INSÉRER]}\\[0.4em]
      \texttt{figures/sprint4\_sequence\_notif.png}\\[0.3em]
      \textit{Séquence notification : Action → Event broadcast → Laravel Reverb →
      Canal privé \texttt{private-project.\{id\}} → Echo.js Frontend → Toast}
    \vspace{1.8cm}}}%
  }
  \caption{Diagramme de séquence — Pipeline de notification WebSocket}
  \label{fig:seq_notif}
\end{figure}

\begin{figure}[H]
  \centering
  \IfFileExists{figures/sprint4_class.png}{%
    \includegraphics[width=0.88\textwidth]{figures/sprint4_class.png}%
  }{%
    \fbox{\parbox{0.88\textwidth}{\centering\vspace{1.8cm}
      \textbf{[IMAGE À INSÉRER]}\\[0.4em]
      \texttt{figures/sprint4\_class.png}\\[0.3em]
      \textit{Classes : Notification, NotificationController,
      MemberDashboardService, AdminDashboardService}
    \vspace{1.8cm}}}%
  }
  \caption{Diagramme de classes — Sprint 4}
  \label{fig:class_s4}
\end{figure}

\subsection{Réalisation}

\subsubsection*{Système de notifications temps réel (Laravel Reverb \& WebSocket)}

Le système de notifications utilise \textbf{Laravel Reverb} (serveur WebSocket natif)
avec \textbf{Laravel Echo} côté frontend. Onze événements sont broadcastés sur des
canaux privés :

\begin{table}[H]
  \centering
  \caption{Événements WebSocket — 11 canaux broadcastés}
  \label{tab:ws_events}
  \begin{tabularx}{\textwidth}{L{4.5cm}X}
    \toprule
    \thead{Événement (classe PHP)} & \thead{Déclencheur} \\
    \midrule
    \rowcolor{lightgray}
    \texttt{TaskUpdated}           & Modification d'une tâche \\
    \texttt{TaskDeleted}           & Suppression de tâche \\
    \rowcolor{lightgray}
    \texttt{TaskAssignmentChanged} & Changement d'assignation \\
    \texttt{TaskCommentCreated}    & Nouveau commentaire \\
    \rowcolor{lightgray}
    \texttt{TaskCommentDeleted}    & Suppression de commentaire \\
    \texttt{ProjectStatusChanged}  & Statut de projet changé \\
    \rowcolor{lightgray}
    \texttt{SprintStateChanged}    & Démarrage / clôture de sprint \\
    \texttt{ProgressUpdated}       & Mise à jour de progression \\
    \rowcolor{lightgray}
    \texttt{NotificationCreated}   & Nouvelle notification in-app \\
    \texttt{NotificationUpdated}   & Notification mise à jour \\
    \rowcolor{lightgray}
    \texttt{ActivityLogged}        & Nouvelle entrée journal d'audit \\
    \bottomrule
  \end{tabularx}
\end{table}

\begin{figure}[H]
  \centering
  \IfFileExists{figures/sprint4_01_notifications.png}{%
    \includegraphics[width=0.85\textwidth]{figures/sprint4_01_notifications.png}%
  }{%
    \fbox{\parbox{0.85\textwidth}{\centering\vspace{2cm}
      \textbf{[IMAGE À INSÉRER]}\\[0.4em]
      \texttt{figures/sprint4\_01\_notifications.png}\\[0.3em]
      \textit{Dropdown notifications : liste avec icône type, timestamp relatif,
      badge compteur non lues, action marquer comme lu}
    \vspace{2cm}}}%
  }
  \caption{Dropdown des notifications temps réel}
  \label{fig:ui_notifs}
\end{figure}

\subsubsection*{Dashboards multi-rôles}

Trois vues de dashboard différenciées sont exposées :
\begin{itemize}
  \item \textbf{Admin} (\texttt{AdminDashboardController}) : métriques globales,
        utilisateurs actifs, projets en cours, accès rapide audit ;
  \item \textbf{Member} (\texttt{MemberDashboardController}) : mes tâches actives,
        sprints en cours, vélocité personnelle ;
  \item \textbf{Client} (\texttt{ClientDashboardController}) : rapport d'avancement
        lecture seule des projets auxquels il est invité.
\end{itemize}

\begin{figure}[H]
  \centering
  \IfFileExists{figures/sprint4_02_dashboard.png}{%
    \includegraphics[width=0.88\textwidth]{figures/sprint4_02_dashboard.png}%
  }{%
    \fbox{\parbox{0.88\textwidth}{\centering\vspace{2cm}
      \textbf{[IMAGE À INSÉRER]}\\[0.4em]
      \texttt{figures/sprint4\_02\_dashboard.png}\\[0.3em]
      \textit{Dashboard Member : KPIs, tâches actives, sprints en cours,
      graphique d'activité personnelle}
    \vspace{2cm}}}%
  }
  \caption{Dashboard personnel (vue Member)}
  \label{fig:ui_dashboard}
\end{figure}

% ============================================================
\section{Sprint 5 : Intelligence artificielle}
% ============================================================

\subsection{Estimation de durée}

L'estimation de durée d'une tâche est réalisée par le service
\texttt{AIService::estimateTaskDuration()} selon une approche heuristique en
plusieurs étapes :
\begin{enumerate}
  \item Calcul d'une durée de base à partir des \textbf{story points}
        (multipliés par \texttt{Task::DEFAULT\_HOURS\_PER\_STORY\_POINT}) ;
  \item Ajustement selon la \textbf{priorité} (urgent: $\times$0.75,
        high: $\times$0.88, low: $\times$1.18) ;
  \item Prise en compte de la \textbf{complexité de la description} (si > 100
        caractères : $\times$1.12) ;
  \item Moyenne pondérée avec la \textbf{durée réelle des tâches similaires}
        passées du même projet.
\end{enumerate}
Le résultat (en minutes) est exposé via \texttt{GET /ai/estimate/\{task\}}.

\subsection{Détection de risques}

Le service \texttt{AIService::detectTaskRisk()} calcule un score de risque
de 0 à 100 par accumulation de points pondérés selon quatre facteurs :

\begin{table}[H]
  \centering
  \caption{Facteurs de calcul du score de risque}
  \label{tab:risk_factors}
  \begin{tabularx}{\textwidth}{L{4cm}X}
    \toprule
    \thead{Facteur} & \thead{Contribution au score} \\
    \midrule
    \rowcolor{lightgray}
    Priorité de la tâche      & Urgent = +30, High = +15, Low = 0 \\
    Proximité de l'échéance   & De 0 à +30 selon les jours restants jusqu'à \texttt{due\_date} \\
    \rowcolor{lightgray}
    Dérapage calendaire       & Si \texttt{planned\_end} dépassé : +20 à +35 \\
    Dépendances bloquantes    & Tâches amont non terminées : +10 à +20 \\
    \bottomrule
  \end{tabularx}
\end{table}

Trois niveaux de risque sont retournés : \texttt{low} (< 30), \texttt{medium}
(30--69) et \texttt{high} ($\geq$ 70). Deux endpoints existent :
\texttt{GET /ai/delay-risk/task/\{task\}} et
\texttt{GET /ai/delay-risk/project/\{project\}}.

\subsection{Recommandation d'affectation}

La recommandation d'assignation repose sur le \textbf{microservice ML FastAPI}
(\texttt{python/ml/app.py}) via le client \texttt{PythonMLClient.php}. Le modèle
(GBM et Random Forest ensemblés via Scikit-learn) est entraîné à partir des
historiques d'assignations de la base de données.

\begin{table}[H]
  \centering
  \caption{Schéma d'entrée du microservice ML (\texttt{POST /recommend})}
  \label{tab:ml_input}
  \begin{tabularx}{\textwidth}{L{3.5cm}L{2.5cm}X}
    \toprule
    \thead{Variable} & \thead{Type} & \thead{Description} \\
    \midrule
    \rowcolor{lightgray}
    \texttt{project\_id}      & int        & Identifiant du projet \\
    \texttt{priority}         & string     & low / medium / high / urgent \\
    \rowcolor{lightgray}
    \texttt{issue\_type}      & string     & task / bug / story \\
    \texttt{complexity}       & int (1–10) & Niveau de complexité estimé \\
    \rowcolor{lightgray}
    \texttt{story\_points}    & int?       & Points si estimés \\
    \texttt{estimated\_hours} & float?     & Heures si estimées \\
    \rowcolor{lightgray}
    \texttt{sprint\_id}       & int?       & Sprint cible \\
    \texttt{specialty\_ids}   & int[]      & Spécialités requises \\
    \rowcolor{lightgray}
    \texttt{top\_n}           & int (1–10) & Nombre de candidats retournés \\
    \bottomrule
  \end{tabularx}
\end{table}

Le modèle retourne un classement des membres par score de probabilité d'assignation
réussie. En cas de données insuffisantes (< 10 échantillons), un modèle de repli
(\textit{LogisticRegression}) est utilisé.

\subsection{Analytics IA}

En complément des fonctionnalités d'estimation et de risque, plusieurs endpoints
d'analyse IA sont disponibles :
\begin{itemize}
  \item \texttt{GET /ai/workload-balance/\{project\}} — analyse de la charge de
        travail par membre ;
  \item \texttt{GET /ai/optimize-assignments/\{project\}} — suggestions d'optimisation
        des assignations à l'échelle du projet ;
  \item \texttt{GET /ai/recommendations/project/\{project\}} — recommandations
        globales sur l'avancement du projet ;
  \item \texttt{GET /projects/\{project\}/auto-assign} — auto-assignation basée
        sur les spécialités (\texttt{SpecialtyController}).
\end{itemize}

\subsection{Spécifications fonctionnelles}

\begin{table}[H]
  \centering
  \caption{Sprint Backlog — Sprint 5 (42 story points)}
  \label{tab:backlog_s5}
  \begin{tabularx}{\textwidth}{C{1.2cm}L{5.8cm}L{2.2cm}C{2cm}C{1.5cm}}
    \toprule
    \thead{ID} & \thead{User Story} & \thead{Acteur} & \thead{Priorité} & \thead{Points} \\
    \midrule
    \rowcolor{lightgray}
    US-24 & Obtenir une recommandation ML d'assignation & Member & Must   & 13 \\
    US-25 & Estimer la durée d'une tâche               & Member & Must   &  8 \\
    \rowcolor{lightgray}
    US-26 & Détecter les risques de retard             & Member & Must   &  8 \\
    US-27 & Discuter avec l'assistant IA               & Member & Should &  8 \\
    \rowcolor{lightgray}
    US-28 & Relancer l'entraînement du modèle ML       & Admin  & Could  &  5 \\
    \midrule
    \multicolumn{4}{r}{\textbf{Total :}} & \textbf{42} \\
    \bottomrule
  \end{tabularx}
\end{table}

\subsection{Conception}

\begin{figure}[H]
  \centering
  \IfFileExists{figures/sprint5_use_case.png}{%
    \includegraphics[width=0.88\textwidth]{figures/sprint5_use_case.png}%
  }{%
    \fbox{\parbox{0.88\textwidth}{\centering\vspace{1.8cm}
      \textbf{[IMAGE À INSÉRER]}\\[0.4em]
      \texttt{figures/sprint5\_use\_case.png}\\[0.3em]
      \textit{Cas d'utilisation Sprint 5 : recommandation ML, estimation, risques,
      chat LLM, réentraînement admin}
    \vspace{1.8cm}}}%
  }
  \caption{Diagramme de cas d'utilisation — Sprint 5}
  \label{fig:uc_s5}
\end{figure}

\begin{figure}[H]
  \centering
  \IfFileExists{figures/sprint5_sequence_ml.png}{%
    \includegraphics[width=0.92\textwidth]{figures/sprint5_sequence_ml.png}%
  }{%
    \fbox{\parbox{0.92\textwidth}{\centering\vspace{1.8cm}
      \textbf{[IMAGE À INSÉRER]}\\[0.4em]
      \texttt{figures/sprint5\_sequence\_ml.png}\\[0.3em]
      \textit{Séquence recommandation : Frontend → GET /ml/recommend/\{task\} →
      MLController → PythonMLClient → FastAPI /recommend →
      GBM+RF modèle → Top-N membres avec scores}
    \vspace{1.8cm}}}%
  }
  \caption{Diagramme de séquence — Pipeline de recommandation ML}
  \label{fig:seq_ml}
\end{figure}

\begin{figure}[H]
  \centering
  \IfFileExists{figures/sprint5_class.png}{%
    \includegraphics[width=0.88\textwidth]{figures/sprint5_class.png}%
  }{%
    \fbox{\parbox{0.88\textwidth}{\centering\vspace{1.8cm}
      \textbf{[IMAGE À INSÉRER]}\\[0.4em]
      \texttt{figures/sprint5\_class.png}\\[0.3em]
      \textit{Classes : AIService, AiController, MLController, PythonMLClient,
      OpenRouterService, AiRecommendation}
    \vspace{1.8cm}}}%
  }
  \caption{Diagramme de classes — Sprint 5}
  \label{fig:class_s5}
\end{figure}

\subsection{Réalisation}

\subsubsection*{Interface de recommandation ML}

\begin{figure}[H]
  \centering
  \IfFileExists{figures/sprint5_01_ml_recommend.png}{%
    \includegraphics[width=0.88\textwidth]{figures/sprint5_01_ml_recommend.png}%
  }{%
    \fbox{\parbox{0.88\textwidth}{\centering\vspace{2cm}
      \textbf{[IMAGE À INSÉRER]}\\[0.4em]
      \texttt{figures/sprint5\_01\_ml\_recommend.png}\\[0.3em]
      \textit{Modal recommandation : liste des membres candidats avec score de
      confiance, badges spécialités, bouton assigner}
    \vspace{2cm}}}%
  }
  \caption{Modal de recommandation ML d'assignation (\texttt{MLRecommendationModal.jsx})}
  \label{fig:ui_ml_recommend}
\end{figure}

\subsubsection*{Assistant conversationnel LLM}

L'assistant (\texttt{AiController::chat}) charge dynamiquement le contexte de
tous les projets accessibles à l'utilisateur via le scope Eloquent
\texttt{Project::accessibleToUser(\$user)} (1 seule requête SQL avec eager loading).
Ce contexte est inclus dans le prompt système envoyé à \textbf{OpenRouter (Llama
3.3 70B Instruct)}.

\begin{figure}[H]
  \centering
  \IfFileExists{figures/sprint5_02_ai_chat.png}{%
    \includegraphics[width=0.85\textwidth]{figures/sprint5_02_ai_chat.png}%
  }{%
    \fbox{\parbox{0.85\textwidth}{\centering\vspace{2cm}
      \textbf{[IMAGE À INSÉRER]}\\[0.4em]
      \texttt{figures/sprint5\_02\_ai\_chat.png}\\[0.3em]
      \textit{Interface chat IA : historique de messages, champ de saisie,
      indicateur de chargement, réponses formatées en Markdown}
    \vspace{2cm}}}%
  }
  \caption{Assistant conversationnel IA — Page \texttt{AiPage.jsx}}
  \label{fig:ui_ai_chat}
\end{figure}

\subsubsection*{Réentraînement du modèle ML}

L'administrateur peut déclencher manuellement le réentraînement via
\texttt{POST /ml/retrain}. Le \texttt{MLController} transmet la requête
à FastAPI \texttt{POST /train}, qui exécute \texttt{recommender.train()} sur les
données de la base. Un minimum de 40 échantillons est requis pour l'entraînement
du modèle ensemblé (GBM + RF), ou 10 pour le modèle de repli (LogisticRegression).

% ============================================================
\section{Conclusion}
% ============================================================

\begin{chapconclusion}
Les sprints~4 et~5 ont apporté les fonctionnalités différenciantes de TAC-TIC
IntelliPM. Le Sprint~4 a livré un système de notifications temps réel couvrant
11 événements WebSocket, des dashboards multi-rôles et la collaboration étendue.
Le Sprint~5 a intégré un moteur IA hybride : recommandation ML via microservice
FastAPI, estimation de durée par heuristiques, détection de risques par règles
métier et assistant conversationnel LLM avec contexte projet dynamique. Le chapitre
suivant présente les modules d'administration enterprise et d'analytique.
\end{chapconclusion}
